\chapter{Lie Superalgebras and Restricted Lie Algebras}

The two theories in this chapter modify different parts of the ordinary
Lie algebra formalism. A superalgebra changes the interchange law for
homogeneous elements. In positive characteristic an ordinary Lie
algebra can carry a \(p\)-operation that is compatible with the
\(p\)-th power in an associative algebra. We prove the corresponding
basis theorems for both enveloping algebras and use them in small
rank examples.

\section{Parity and super enveloping algebras}

\begin{definition}[Lie superalgebras and their modules]
\label{res:SR01}
Let \(k\) be a field of characteristic different from two.
A Lie superalgebra is a graded vector space
\(\mathfrak a=\mathfrak a_{\overline 0}\oplus\mathfrak a_{\overline 1}\)
with an even bilinear bracket, meaning that degrees add. For
homogeneous \(x,y,z\), write
\(\epsilon_{xy}=(-1)^{|x||y|}\).
The identities are
\[
 [x,y]=-\epsilon_{xy}[y,x],
 \qquad [x,[y,z]]=[[x,y],z]+\epsilon_{xy}[y,[x,z]].
\]
In characteristic three we also require
\([x,[x,x]]=0\) for each odd \(x\). In other allowed characteristics
this follows from the first two identities. A graded module is a
graded vector space \(M\) with maps \(x\mapsto \rho(x)\) carrying
parity \(|x|\), satisfying
\(\rho([x,y])=\rho(x)\rho(y)-\epsilon_{xy}\rho(y)\rho(x)\).
Morphisms preserve the grading.
\end{definition}

Throughout the examples involving complex representations, the
ground field is \(\mathbb C\). The parity and basis arguments
use only the hypotheses of Definition~\ref{res:SR01}.
For the superalgebra conventions see Musson
~\cite[Section 1.1, PDF pp.~22--24]{musson}
and Serganova~\cite[Section 2]{serganova}.

\begin{proposition}[The supercommutator and the smallest matrix example]
\label{res:SR02}
In an associative graded algebra the bracket
\([x,y]=xy-\epsilon_{xy}yx\) makes its graded vector space a
Lie superalgebra. In \(\mathfrak{gl}(1|1)\), with even
generators \(z=E_{11}+E_{22}\), \(h=E_{11}\) and odd
generators \(e=E_{12}\), \(f=E_{21}\), the nonzero brackets
up to graded skew symmetry are
\[
 [h,e]=e,\qquad [h,f]=-f,\qquad [e,f]=z .
\]
In particular \(z\) is central and \(e,f\) are odd with
zero self-bracket. The supertrace is a linear functional
on \(\mathfrak{gl}(1|1)\) vanishing on every bracket, with
kernel \(\operatorname{span}_k(z,e,f)\).
\end{proposition}

\begin{proof}
Skew symmetry follows by interchanging the two products.
Expanding the derivation identity for \([x,-]\) on a
product gives
\[
 [x,yz]=[x,y]z+\epsilon_{xy}y[x,z]:
\]
the intermediate \(yxz\) terms cancel. Apply this to both
products in \([y,z]\) and use
\(\epsilon_{x,yz}=\epsilon_{xy}\epsilon_{xz}\) to obtain
the graded Jacobi identity. If \(x\) is odd then
\([x,x]=2x^2\), which commutes with \(x\). This also
verifies the extra identity in characteristic three.
For an abstract odd \(x\), substituting \(x=y=z\) into
Jacobi and using skew symmetry gives \(3[x,[x,x]]=0\).

For matrices take block-diagonal elements to be even
and off-diagonal elements odd. Multiplication of the
matrix units gives the displayed brackets
and \(e^2=f^2=0\).
The supertrace sends \(\operatorname{diag}(a,b)\)
to \(a-b\) and sends odd elements to zero. The
bracket table shows directly that its kernel is
the asserted span and contains the bracket subspace.
\end{proof}

\begin{theorem}[Super Poincare--Birkhoff--Witt]
\label{res:SR03}
Let \(\mathfrak a\) satisfy Definition~\ref{res:SR01}, and
order a homogeneous basis. Its enveloping superalgebra is
the quotient of \(T_k(\mathfrak a)\) by
\[
 xy-\epsilon_{xy}yx-[x,y]\qquad(x,y\text{ homogeneous}).
\]
The weakly increasing words in this basis, with each
odd basis element appearing at most once, form a basis.
For the word-length filtration its associated graded is
\[
 \operatorname{gr} U(\mathfrak a)
   \cong S(\mathfrak a_{\overline 0})
              \otimes_k\bigwedge(\mathfrak a_{\overline 1})
\]
as graded superalgebras. Graded modules over the enveloping
algebra and graded Lie superalgebra modules agree.
\end{theorem}

\begin{proof}
For a descending pair \(a>b\) replace \(ab\) by
\(\epsilon_{ab}ba+[a,b]\). For an odd repeated letter
replace \(aa\) by \([a,a]/2\). Order reductions first
by length and then by number of inversions. Each
replacement has smaller terms for this order, so
recursive reduction terminates. The relations of the
presentation are generated by these replacements:
for a weakly increasing pair use graded skew symmetry
or the repeated-letter relation, and extend bilinearly.

Disjoint adjacent replacements commute. Two overlapping
replacements have input \(abc\) with \(a\ge b\ge c\),
where a repeated adjacent pair must be odd. If all
three are distinct, reorder the three-letter term to
\(\epsilon_{ab}\epsilon_{ac}\epsilon_{bc}cba\) in
both paths. The difference of the shorter terms reduces
to
\[
 [[a,b],c]+\epsilon_{ab}[b,[a,c]]-[a,[b,c]]=0.
\]
Indeed the respective pairs of shorter terms are
\([a,b]c-\epsilon_{ac}\epsilon_{bc}c[a,b]\),
\(\epsilon_{ab}b[a,c]-\epsilon_{bc}[a,c]b\), and
\(\epsilon_{ab}\epsilon_{ac}[b,c]a-a[b,c]\).
For input \(aac\) with \(a\) odd, the difference reduces to
\(\tfrac12[[a,a],c]-[a,[a,c]]=0\).
For \(abb\) with \(b\) odd, it reduces to
\([[a,b],b]-\tfrac12[a,[b,b]]=0\).
These equalities are Jacobi with repeated odd entries.
Finally for \(aaa\) with \(a\) odd the difference reduces to
\(\tfrac12[[a,a],a]=0\), using the odd triple identity.
The checks are valid for bracket values that are sums of
basis elements by linearity.

Induct on the length-inversion order to compare any two possible
first steps. In the disjoint case the two paths have the same
common expansion. In the overlap case the expansion comparison
above leaves terms of smaller length, whose reductions agree by
induction. Prefixes and suffixes preserve this smaller-length
comparison. This gives a unique linear normal form. It kills
each replacement relation with every prefix and suffix by
the just-proved independence of first step, and so kills
their two-sided ideal. Conversely each reduction differs
from its input by an element of that ideal. Thus the normal
words are a basis of the quotient. Their leading symbols
give the displayed symmetric and exterior factors:
the leading relations say that even letters commute and
odd letters anticommute and square to zero. The tensor
algebra action on a graded module descends exactly when it
satisfies the signed bracket identity in
Definition~\ref{res:SR01}, giving the module claim.
\end{proof}

\begin{example}[Simple graded modules for the complex \(\mathfrak{gl}(1|1)\)]
\label{res:SR04}
On a finite-dimensional simple graded module the central \(z\)
acts as a scalar \(\gamma\). If \(\gamma=0\), the module is
one-dimensional, with \(e=f=0\) and \(h\) an arbitrary scalar
\(\lambda\). If \(\gamma\ne0\), it has a homogeneous basis
\(v,w\) of opposite parities on which
\[
\begin{gathered}
 hv=\lambda v,\qquad hw=(\lambda-1)w,\qquad
 fv=w,\qquad fw=0,\\
 ev=0,\qquad ew=\gamma v,\qquad
 zv=\gamma v,\qquad zw=\gamma w .
\end{gathered}
\]
In each case the scalar data and the parity of \(v\)
determine the isomorphism class in the graded category.
\end{example}

\begin{proof}
The kernel of \(z-\gamma\) for an eigenvalue \(\gamma\)
is a nonzero graded submodule, hence is the whole module.
Since \(e^2=0\), \(\ker e\ne0\). It is graded and
stable under \(h\), so over \(\mathbb C\) it contains a
homogeneous \(h\)-eigenvector \(v\), of weight \(\lambda\).
The table in Proposition~\ref{res:SR02} gives
\(e(fv)=\gamma v\), \(f^2v=0\), and \(h(fv)=(\lambda-1)fv\).
Thus \(\operatorname{span}(v,fv)\) is the whole simple module.
If \(\gamma=0\) and \(fv\ne0\), the line it spans would be a
nonzero proper graded submodule. If \(\gamma\ne0\), the
vector \(w=fv\) is nonzero. The displayed formulas define
a graded representation by the bracket table. It is simple
because either weight line generates both lines, and a nonzero
invariant subspace has a weight vector by interpolation in \(h\).
The weight and parity of the one-dimensional kernel of \(e\)
distinguish the parameters. In this two-dimensional case the
opposite choice of positive odd root instead uses \(\ker f\),
so changes the highest weight to \(\lambda-1\) and changes
its parity. On the one-dimensional modules both choices
give the same highest line.
\end{proof}

\section{The \texorpdfstring{\(p\)}{p}-operation}

For the rest of the chapter, \(k\) has characteristic \(p>0\),
and all Lie algebras are ordinary. Write \(U(\mathfrak g)\)
for the ordinary universal enveloping algebra. We use the
ordered-basis and associated-graded parts of the ordinary PBW
theorem in \emph{Lie Algebras and Finite-Dimensional
Representations}, Result A05~\cite[PDF pp.~8--9]{mg-lie}.
That proof applies to any characteristic for a Lie algebra
with an ordered \(k\)-basis.

\begin{lemma}[Associative \(p\)-powers and the Jacobson polynomials]
\label{res:SR05}
Let \(B\) be an associative \(k\)-algebra, with ordinary
commutator bracket. For \(1\le i\le p-1\), define
\(s_i(x,y)\) by specifying that \(i\,s_i(x,y)\) is the
coefficient of \(t^{i-1}\) in
\((\operatorname{ad}(t x+y))^{p-1}(x)\).
These are Lie polynomials in \(x,y\). One has
\[
 (\operatorname{ad}x)^p=\operatorname{ad}(x^p),\qquad
 (x+y)^p=x^p+y^p+\sum_{i=1}^{p-1}s_i(x,y).
\]
\end{lemma}

\begin{proof}
Left and right multiplication by \(x\) commute.
In characteristic \(p\) the \(p\)-th power of their
difference is the difference of their \(p\)-th powers,
proving the first equality. Set \(X=t x+y\) in \(B[t]\).
Formal differentiation of \(X^p\), using a central variable,
gives
\[
 \frac{d}{dt}X^p=\sum_{j=0}^{p-1}X^{p-1-j}xX^j
          =(\operatorname{ad}X)^{p-1}(x),
\]
since \((-1)^j\binom{p-1}{j}=1\) in \(k\).
The constant and degree-\(p\) coefficients of \(X^p\)
are \(y^p,x^p\). They contribute zero to the derivative.
For each intermediate coefficient, differentiation
multiplies it by the corresponding invertible integer \(i\).
It is therefore \(s_i(x,y)\). Evaluating at \(t=1\)
gives the second equality. The iterated-commutator
description proves that these coefficients are Lie
polynomials and makes their definition available in any
ordinary Lie algebra.
\end{proof}

\begin{definition}[Restricted Lie algebras]
\label{res:SR06}
A restricted Lie algebra is an ordinary Lie algebra with a map
\(x\mapsto x^{[p]}\) satisfying
\[
\begin{aligned}
\operatorname{ad}(x^{[p]})&=(\operatorname{ad}x)^p,\\
(a x)^{[p]}&=a^p x^{[p]},\\
(x+y)^{[p]}&=x^{[p]}+y^{[p]}+\sum_{i=1}^{p-1}s_i(x,y),
\end{aligned}
\]
with the same polynomials as in Lemma~\ref{res:SR05}.
A Lie subalgebra closed under the operation inherits it.
\end{definition}

\begin{proposition}[Matrix and derivation \(p\)-operations]
\label{res:SR07}
Associative \(p\)-th powers restrict to \(\mathfrak{sl}_n(k)\)
inside the full matrix algebra and give a restricted
structure there. They also give a restricted structure on
\(\operatorname{Der}_k(A)\) for a commutative \(k\)-algebra \(A\),
using composition in \(\operatorname{End}_k(A)\).
In particular for \(A=k[t]/(t^p)\) and
\(D_i=t^{i+1}\frac{d}{dt}\), \(-1\le i\le p-2\), one has
\[
 D_0^{[p]}=D_0,\qquad D_i^{[p]}=0\quad(i\ne0).
\]
\end{proposition}

\begin{proof}
Over an algebraic closure a matrix is triangularizable by
inductively choosing an eigenline and passing to the quotient.
Its \(p\)-th power has the \(p\)-th powers of its eigenvalues
on the diagonal. Thus
\(\operatorname{tr}(X^p)=\operatorname{tr}(X)^p\), an identity
which holds already over \(k\). The traceless Lie subalgebra
is therefore closed under the power. For a derivation \(D\),
repeated Leibniz gives
\[
 D^p(ab)=\sum_{j=0}^p\binom pj D^j(a)D^{p-j}(b)
         =D^p(a)b+aD^p(b).
\]
In both cases the restricted identities are those of
Lemma~\ref{res:SR05} and scaling.

Every derivation of \(k[t]/(t^p)\) is determined by the
image of \(t\), which can be arbitrary since the derivative
of \(t^p\) is zero. The displayed \(D_i\) form a basis.
The operator \(D_0\) acts on \(t^m\) by the scalar
\(m\in\mathbb F_p\), whose \(p\)-th power is itself.
On monomials \(D_{-1}\) lowers their degree, and \(D_i\)
with \(i>0\) raises it until the monomial vanishes in
the quotient. Their \(p\)-fold compositions are zero.
For \(p>3\), the simplicity and vanishing Killing form
of this Witt algebra were proved in the earlier volume,
Result J09~\cite[PDF p.~72]{mg-lie}.
\end{proof}

\section{The \texorpdfstring{\(p\)}{p}-center and reduced enveloping algebras}

\begin{theorem}[Polynomial \(p\)-center and reduced PBW]
\label{res:SR08}
Let \(\mathfrak g\) be a restricted Lie algebra of dimension
\(n<\infty\). The function
\[
 z(x)=x^p-x^{[p]}\in U(\mathfrak g)
\]
is central, additive and \(p\)-semilinear:
\(z(a x)=a^p z(x)\). If \(x_1,\ldots,x_n\) is a basis,
the subalgebra \(Z_p\) generated by all \(z(x)\) is the
polynomial ring on \(z(x_1),\ldots,z(x_n)\).
The algebra \(U(\mathfrak g)\) is free over \(Z_p\) on the
ordered monomials with each exponent less than \(p\).

For a linear functional \(\chi:\mathfrak g\to k\), set
\[
 u_\chi(\mathfrak g)
   = U(\mathfrak g)\big/\bigl(z(x)-\chi(x)^p:x\in\mathfrak g\bigr).
\]
The images of those same ordered monomials are a
\(k\)-basis, so the dimension is \(p^n\). In particular
\(u_0(\mathfrak g)\) is the restricted enveloping algebra.
\end{theorem}

\begin{proof}
For \(y\in\mathfrak g\), Lemma~\ref{res:SR05} and the
restricted adjoint identity give
\([x^p,y]=(\operatorname{ad}x)^p(y)=[x^{[p]},y]\);
the \(y\)'s generate \(U(\mathfrak g)\). The two
Jacobson sum formulas for associative powers and the
restricted map cancel upon subtraction, proving additivity.
The scaling identity is immediate. Thus the \(z(x_i)\)
already generate \(Z_p\).

In the associated graded algebra of ordinary PBW,
the leading symbol of \(z(x_i)\) is \(x_i^p\), since
the subtracted term has filtration degree at most one.
Write \(a_i=p b_i+r_i\) with \(0\le r_i<p\) for any
ordinary PBW word of exponents \((a_i)\).
The element
\(\prod_i z(x_i)^{b_i}\prod_i x_i^{r_i}\),
with the \(x_i\) in their fixed order, has that word's
leading symbol. Subtract and use induction on the degree
to express every element as a linear combination of these
products. For independence, distinct pairs \((b,r)\)
have distinct associated-graded leading monomials.
In a finite relation the terms of maximal degree cannot
cancel there. This proves polynomial independence and
the freeness assertion.

By \(p\)-semilinearity the ideal in the definition of
\(u_\chi\) is generated by the central elements
\(z(x_i)-\chi(x_i)^p\). On the free \(Z_p\)-module just
constructed, quotienting by that ideal specializes each
polynomial variable to \(\chi(x_i)^p\). This leaves exactly
the stated basis over \(k\).
\end{proof}

\begin{corollary}[The \(p\)-character of a finite-dimensional simple module]
\label{res:SR09}
If \(k\) is algebraically closed, every finite-dimensional simple
\(\mathfrak g\)-module has a unique linear functional \(\chi\)
such that it is a module over \(u_\chi(\mathfrak g)\).
Conversely a simple module of any \(u_\chi(\mathfrak g)\) is
a finite-dimensional simple \(\mathfrak g\)-module with this
\(p\)-character.
\end{corollary}

\begin{proof}
On a finite-dimensional simple, each central \(z(x_i)\)
acts as a scalar \(c_i\): a nonzero eigenspace is a
submodule. The \(p\)-power map on an algebraically closed
field of characteristic \(p\) is bijective, since it
is injective on a field and every \(T^p-c_i\) has a root.
There is a unique linear \(\chi\) with
\(\chi(x_i)^p=c_i\), and Lemma~\ref{res:SR05} and
Theorem~\ref{res:SR08} ensure the remaining relations
by semilinearity. Conversely \(u_\chi\) is finite-dimensional
and a simple module is generated by any one nonzero vector,
so is finite-dimensional. Submodules for the quotient are
the same as invariant subspaces for the Lie generators,
proving simplicity and the assertion about \(\chi\).
\end{proof}

For the restricted identities and reduced algebra see
Jantzen~\cite[Sections A.1--A.3 and A.7, PDF pp.~1--3 and 4--5]{jantzen-modular}
and Strade--Farnsteiner~\cite[Sections 2.1, 2.5 and 5.3]{stradefarnsteiner}.

\section{Baby Verma modules for \texorpdfstring{\(\mathfrak{sl}_2\)}{sl(2)}}

Assume now that \(k\) is algebraically closed and \(p>2\).
Put \(e=E_{12}\), \(f=E_{21}\), and \(h=E_{11}-E_{22}\).
The brackets are \([h,e]=2e\), \([h,f]=-2f\), and
\([e,f]=h\). By Proposition~\ref{res:SR07} one has
\(e^{[p]}=f^{[p]}=0\) and \(h^{[p]}=h\).
This section computes the simple reduced modules in all
three types of \(p\)-character, separating their highest
weight labels from the characteristic-zero labels.

\begin{proposition}[Baby Verma induction]
\label{res:SR10}
Suppose \(\chi(e)=0\). Write \(\eta=\chi(h)\) and
\(\xi=\chi(f)\). For each root \(\lambda\) of
\(\lambda^p-\lambda=\eta^p\), there is a module \(Z_\chi(\lambda)\)
with basis \(v_i=f^i v_0\), \(0\le i<p\), and
\[
\begin{aligned}
 hv_i&=(\lambda-2i)v_i,\\
 ev_0&=0,\qquad ev_i=i(\lambda-i+1)v_{i-1}\quad(1\le i<p),\\
 fv_i&=v_{i+1}\quad(i<p-1),\qquad fv_{p-1}=\xi^p v_0.
\end{aligned}
\]
Every simple \(u_\chi(\mathfrak{sl}_2)\)-module is a
quotient of one of these modules.
\end{proposition}

\begin{proof}
For \(\mathfrak b=k h+k e\), the one-dimensional space
with \(h\) acting by \(\lambda\) and \(e\) by zero respects
the bracket and the reduced power identities, so is a
\(u_{\chi|_{\mathfrak b}}(\mathfrak b)\)-module. Here
\(\mathfrak b\) is closed under the \(p\)-map because powers
of an upper triangular matrix stay upper triangular.
Order the reduced PBW monomials as \(f^i h^j e^\ell\).
The basis in Theorem~\ref{res:SR08} identifies the reduced
algebra of \(\mathfrak b\) with its subalgebra in the
reduced algebra of \(\mathfrak{sl}_2\), and shows the
latter is free as a right module over it on \(f^i\).
Inducing the one-dimensional module therefore gives a
\(p\)-dimensional module with the asserted basis.
The \(h\) and \(f\) formulas follow from the bracket
and the defining reduced relations. Commuting \(e\)
past one further \(f\), the coefficient in
\(e f^{i+1}v_0\) is
\((\lambda-2i)+i(\lambda-i+1)=(i+1)(\lambda-i)\).
This proves the \(e\) formula by induction.

On any simple \(u_\chi(\mathfrak{sl}_2)\)-module \(e^p=0\), so
\(\ker e\ne0\), and \([h,e]=2e\) implies this kernel
is \(h\)-stable. Choose an \(h\)-eigenvector in it.
Its weight satisfies the reduced relation for \(h\).
The induced module with that weight maps onto the simple,
because the nonzero vector generates it.
\end{proof}

\begin{theorem}[The zero \(p\)-character]
\label{res:SR11}
For \(\chi=0\) the simple modules are \(L(m)\) for
\(m=0,\ldots,p-1\). They have dimensions \(m+1\).
For \(\lambda=m\) in the prime field, \(L(m)\) is the
quotient of \(Z_0(m)\) by the span of
\(v_{m+1},\ldots,v_{p-1}\), with empty span for \(m=p-1\).
\end{theorem}

\begin{proof}
The solutions of \(\lambda^p-\lambda=0\) are the prime
field elements. The \(p\) weights of \(Z_0(m)\) are
distinct because \(2\ne0\) in \(k\). The displayed tail
is stable: its first vector is killed by \(e\), and \(f\)
kills \(v_{p-1}\). In the quotient a nonzero submodule
contains a weight vector \(v_i\) with \(0\le i\le m\),
since projections onto the finitely many weight spaces
are polynomials in \(h\). For this range every noninitial
\(e\)-coefficient is nonzero. Repeated application of \(e\)
reaches \(v_0\) and \(f\) then produces the entire quotient.
This proves simplicity. The same argument in \(Z_0(m)\)
shows that any submodule containing a vector with nonzero
component in one of the first \(m+1\) weight spaces is
the whole module. Thus its proper submodules all lie in
the tail and the simple quotient is unique.
The modules \(L(m)\) are pairwise nonisomorphic by their
distinct dimensions, and Proposition~\ref{res:SR10}
gives exhaustion.
\end{proof}

\begin{theorem}[Nonzero semisimple and nilpotent \(p\)-characters]
\label{res:SR12}
There are the following two nonzero cases, after a change
of basis by conjugation of matrices.

\emph{Semisimple case.} Suppose
\(\chi(e)=\chi(f)=0\), \(\eta=\chi(h)\ne0\).
The \(p\) modules \(Z_\chi(\lambda)\) for the roots of
\(\lambda^p-\lambda=\eta^p\) are simple and pairwise
nonisomorphic, each of dimension \(p\).

\emph{Nilpotent case.} Suppose
\(\chi(e)=\chi(h)=0\), \(\chi(f)=1\).
All \(p\) baby Verma modules \(Z_\chi(m)\), \(m\in\mathbb F_p\),
are simple of dimension \(p\). The isomorphism classes
are indexed by the orbits of \(m\mapsto -m-2\), giving
\((p+1)/2\) classes.
\end{theorem}

\begin{proof}
In the semisimple case the polynomial of the weights is
separable of degree \(p\) and has no root in the prime field.
All \(e\)-coefficients for \(1\le i<p\) in
Proposition~\ref{res:SR10} are nonzero. A nonzero
submodule contains a weight line, and repeated applications
of \(e\) and \(f\) generate the entire basis. The kernel of
\(e\) is only the highest line, whose weight distinguishes
the \(p\) modules. They exhaust by induction.

In the nilpotent case the \(h\)-weights are still distinct
and \(f\) cyclically permutes the \(p\) weight lines since
\(f^p=1\). A submodule containing one weight line therefore
contains them all. For \(m\ne-1\) the kernel of \(e\) is
spanned by \(v_0,v_{m+1}\), where \(m\) is represented by
its integer in \(0,\ldots,p-2\). Their weights are \(m\)
and \(m-2(m+1)=-m-2\). Each is a highest line and induces
a map onto the simple, so the modules with these two labels
are isomorphic. For \(m=-1\) only \(v_0\) is killed by \(e\).
Every isomorphism preserves the weights in \(\ker e\), so
there are no other identifications. There is exactly one
fixed element of the involution, \(-1\), giving the count.

Finally the trace pairing on \(\mathfrak{sl}_2(k)\) is
nondegenerate: it pairs \(e\) with \(f\) and has
\(\operatorname{tr}(h^2)=2\). Every \(\chi\) is therefore
\(\operatorname{tr}(X\,\cdot)\) for one traceless \(X\).
A nonzero such matrix is conjugate either to
\(\operatorname{diag}(a,-a)\) with \(a\ne0\) or to
\(E_{12}\), according as its two eigenvalues are distinct
or both zero. The displayed parameters follow from the
pairing. Conjugation preserves the bracket and matrix
\(p\)-power, so it transports the corresponding reduced
algebras and their modules. This proves the scope of the
two cases.
\end{proof}

The inductive construction follows the setting of
Jantzen~\cite[Sections B.3--B.4, B.10 and C.3,
PDF pp.~6--7 and 9--11]{jantzen-modular}.
General classification of simple modular Lie algebras and of
simple modules for higher-rank reduced enveloping algebras
uses further hypotheses and results. Here the module
classifications are for \(\mathfrak{gl}(1|1)\) over
\(\mathbb C\) and for ordinary \(\mathfrak{sl}_2\) over an
algebraically closed field of characteristic greater than two.
